\section{Encoding Rules}
\label{sec:beer-encoding-rules}

This section defines the encoder as a mathematical object: a partial function from typed B terms to typed SMT terms.
The rules below are the monad-free reading of the Lean function {\small\leancodeptr{encodeTerm}{Encoder/Encoder.html\#encodeTerm}}; how they are realized as a stateful program is deferred to \Cref{sec:beer-implementation}.

\subsection{The encoding judgement}
\label{sec:beer-encoding-judgement}

We write
\[
  t \encto (t', \sigma)
\]
when the encoder, applied to the B term \(t\), succeeds and returns the SMT term \(t'\) together with its SMT type \(\sigma\).
The judgement is partial: on unsupported constructs, or on operands whose representations cannot be reconciled, the encoder fails (written \(\throw\), as in \Cref{ch:loosening}) and no SMT term is produced.

Two parameters are left implicit in the notation, both managed by the encoder state (\Cref{sec:beer-implementation}).
First, an ambient SMT typing context \(\Gamma\) assigning an SMT type to every variable in scope; the rules below consult it for variables and extend it under binders.
Second, a supply of fresh SMT variables, drawn upon whenever a rule introduces a binder or a helper constant.
The judgement maintains the invariant that guides the whole design: a term of B type \(\alpha\) encodes to an SMT term of type \(\toSMT{\alpha}\),
\[
  \btype{\Gamma_{\mathsf{B}}}{t}{\alpha} \;\wedge\; t \encto (t', \sigma)
  \quad\Longrightarrow\quad
  \sigma = \toSMT{\alpha} \;\wedge\; \stype{\Gamma}{t'}{\toSMT{\alpha}},
\]
when \(\Gamma\) assigns to each variable the encoding of its B type.
This typing-preservation property is part of the soundness theorem and is proved in \Cref{sec:beer-soundness}; a deliberate exception for \emph{flagged} variables, which are assigned a strictly tighter type, is introduced in \Cref{sec:beer-rules-cast}.

\subsection{Atoms, arithmetic, logic, and pairs}
\label{sec:beer-rules-homomorphic}

The base fragment is homomorphic: each B operator maps to its SMT counterpart, with operands encoded recursively.
\begin{align*}
  \var{v} &\encto (\var{v},\, \Gamma(\var{v}))
  & n &\encto (n,\, \intS)
  & b &\encto (b,\, \boolS) \\
  x \addB y &\encto (x' \addS y',\, \intS)
  & x \subB y &\encto (x' \subS y',\, \intS)
  & x \mulB y &\encto (x' \mulS y',\, \intS) \\
  x \leB y &\encto (x' \leS y',\, \boolS)
  & x \andB y &\encto (x' \andS y',\, \boolS)
  & \notB x &\encto (\notS x',\, \boolS)
\end{align*}
where \(x \encto (x', \cdot)\) and \(y \encto (y', \cdot)\), the arithmetic rules requiring the operands to encode at type \(\intS\) and the logical ones at \(\boolS\).
Variables are translated to themselves---B and SMT variable names are shared---at the type recorded in \(\Gamma\).
The maplet becomes a pair:
\[
  x \mapletB y \encto (\pairS{x'}{y'},\, \sigma_x \timesS \sigma_y)
  \qquad\text{where } x \encto (x', \sigma_x),\; y \encto (y', \sigma_y).
\]
No rules are needed for \(\orB\), \(\impB\), or \(\iffB\): as noted in \Cref{ch:b}, B's connectives are generated by \(\andB\) and \(\notB\), so these forms are eliminated before the encoder ever sees them.

\subsection{Sets as characteristic predicates}
\label{sec:beer-rules-sets}

Following \Cref{def:beer-type-encoding}, a B set of type \(\setB\alpha\) encodes as a term of type \(\toSMT{\alpha} \toS \boolS\).
The two base-type sets become constant-true predicates over the corresponding SMT sort:
\[
  \ZintB \encto (\lamS \var{z} : \intS \cdot \mathsf{true},\; \intS \toS \boolS)
  \qquad
  \ZboolB \encto (\lamS \var{z} : \boolS \cdot \mathsf{true},\; \boolS \toS \boolS)
\]
The set formers each emit a lambda whose body is their defining predicate.
For the powerset, with \(S \encto (S', \sigma \toS \boolS)\),
\[
  \powB(S) \encto \Bigl(\lamS \var{E} : \sigma \toS \boolS \cdot\; \forallS \var{x} : \sigma \cdot \appS{\var{E}}{\var{x}} \impS \appS{S'}{\var{x}},\;\; (\sigma \toS \boolS) \toS \boolS\Bigr),
\]
the predicate expressing \(\var{E} \subseteq S'\) pointwise.
For the Cartesian product, with \(A \encto (A', \sigma_1 \toS \boolS)\) and \(B \encto (B', \sigma_2 \toS \boolS)\),
\[
  A \prodB B \encto \Bigl(\lamS \var{p} : \sigma_1 \timesS \sigma_2 \cdot\; \existsS \var{a}, \var{b} \cdot \appS{A'}{\var{a}} \andS \appS{B'}{\var{b}} \andS (\var{p} \eqS \pairS{\var{a}}{\var{b}}),\;\; (\sigma_1 \timesS \sigma_2) \toS \boolS\Bigr).
\]
For the partial-function space, the emitted predicate is the conjunction of the two defining properties of a partial function---its graph is contained in \(A \prodB B\), and it is functional:
\[
  A \pfunB B \encto \Bigl(\lamS \var{R} \cdot\;
  \bigl(\forallS \var{x}, \var{y} \cdot \appS{\var{R}}{\pairS{\var{x}}{\var{y}}} \impS \appS{A'}{\var{x}} \andS \appS{B'}{\var{y}}\bigr)
  \andS
  \bigl(\forallS \var{x}, \var{y}, \var{y'} \cdot \appS{\var{R}}{\pairS{\var{x}}{\var{y}}} \andS \appS{\var{R}}{\pairS{\var{x}}{\var{y'}}} \impS \var{y} \eqS \var{y'}\bigr),\;\;
  ((\sigma_1 \timesS \sigma_2) \toS \boolS) \toS \boolS\Bigr)
\]
with \(\var{R} : (\sigma_1 \timesS \sigma_2) \toS \boolS\).
All richer B function spaces (\(\tfunB\), \(\injpfunB\), \dots) are comprehensions over \(\pfunB\) (\Cref{ch:b}) and therefore reach the encoder in already-elaborated form.

\subsection{Comprehension and lambda}
\label{sec:beer-rules-comprehension}

A set comprehension over a set-typed domain \(D \encto (D', \sigma \toS \boolS)\) encodes as a guarded characteristic predicate:
\[
  \bcomp{\vec{\var{v}} \inB D}{P} \encto \Bigl(\lamS \var{z} : \sigma \cdot\; \iteS\bigl(\appS{D'}{\var{z}},\; P'[\vec{\var{v}} \leftarrow \mathsf{destr}(\var{z})],\; \mathsf{false}\bigr),\;\; \sigma \toS \boolS\Bigr)
\]
where \(P \encto (P', \boolS)\) is encoded in the context extended with the bound variables, and \(\mathsf{destr}_i(\var{z})\) is the \(i\)-th projection of \(\var{z}\), a chain of \(\fstS\)/\(\sndS\) applications destructuring the tuple variable into the components bound by \(\vec{\var{v}}\).
The guard confines the predicate to \(D'\): outside it, the body---which may be meaningless there---is never consulted, and the predicate returns \(\mathsf{false}\).
This guard is what reconciles the partiality of the B construct with the totality that SMT lambdas demand; the point recurs in the soundness proof (\Cref{sec:beer-case-comprehension}).

A B lambda encodes as the \emph{graph} of the partial function it denotes.
With \(D \encto (D', \sigma \toS \boolS)\) and \(P \encto (P', \gamma)\),
\[
  \lambdaB \vec{\var{v}} \inB D \cdot P \encto \Bigl(\lamS \var{u} : \sigma \timesS \gamma \cdot\; \appS{D'}{(\fstS \var{u})} \andS \bigl(\sndS \var{u} \eqS P'[\vec{\var{v}} \leftarrow \mathsf{destr}(\fstS \var{u})]\bigr),\;\; (\sigma \timesS \gamma) \toS \boolS\Bigr),
\]
the characteristic predicate of \(\{(x, y) \mid x \in D \wedge y = P(x)\}\).
This matches the SMT encoding of the B type: a B lambda has type \(\setB(\tau \prodB \beta)\), a set of pairs, not a function type.
Encoding the graph rather than a function value also spares the encoder any commitment about the lambda's value outside \(D\), where B leaves it undefined.

For both constructs, the implementation contains a second branch for domains encoded at a \emph{functional} type (a flagged domain, \Cref{sec:beer-rules-cast}); it emits an \(\iteS\)/\(\optS\)-based variant, \eg \(\lamS \vec{\var{x}} \cdot \iteS(P', \someS D'(\vec{\var{x}}), \noneS)\) for comprehensions.
These variants lie outside the fragment covered by the soundness proof and are discussed in \Cref{sec:beer-limitations}.

\subsection{The cast layer}
\label{sec:beer-rules-cast}

The rules so far assign every term its canonical type \(\toSMT{\alpha}\).
The encoder, however, deliberately breaks canonicity for one class of variables.
The Atelier~B proof-obligation format marks certain identifiers as \emph{functional}---the decoder flags a variable \(\var{x}\) when it appears in a hypothesis of the form \(\var{x} = \lambdaB\ldots\), \(\var{x} \inB A \pfunB B\), \(\var{x} \inB\) a comprehension over \(\pfunB\), or an equality with an already-flagged variable.
A flagged variable of B type \(\setB(\alpha \prodB \beta)\) is declared at the SMT type
\[
  \toSMT{\alpha} \toS \optS\toSMT{\beta}
  \qquad\text{instead of the canonical}\qquad
  (\toSMT{\alpha} \timesS \toSMT{\beta}) \toS \boolS,
\]
a partial-function representation that solvers exploit far better than a bare graph predicate.
The cost is that one B type now inhabits several \(\sqsubseteq\)-comparable SMT types---exactly the multi-representation problem of \Cref{ch:loosening}---and every binary operation must first bring its operands to a common representation.
This is the task of the cast combinators.
Each follows the same scheme: compare the operand types with the decidable order \(\sqsubseteq\); if they differ, loosen the tighter operand along the cast path to the more general type (their join), obtaining a fresh variable \(\lvar{x}\) constrained by its specification \(\spec{\lvar{x}}\) to denote the semantic cast image of \(x\) (\Cref{thm:loosening-correctness}); then combine pointwise at the common type.
Incomparable types are rejected (\(\throw\)).

\paragraph{Union and intersection.}
When both operands already share a characteristic-predicate type, no cast is needed and the encoder combines them directly ({\small\leancodeptr{castUnion}{Encoder/Loosening/Loosening.html\#castUnion}}): for \(S, T\) of common type \(\gamma \toS \boolS\),
\begin{equation}
  S \cupB T \encto \bigl(\lamS \var{z} : \gamma \cdot\; \appS{S}{\var{z}} \orS \appS{T}{\var{z}},\;\; \gamma \toS \boolS\bigr),
  \label{eq:beer-union-direct}
\end{equation}
and intersection likewise with \(\andS\).
Otherwise the shape of the cast path dictates the rule.
The representative case---exercised by the running example---is the \(\mathsf{graph}\) path: \(S : \alpha \toS \optS\beta\) functional, \(T : (\alpha' \timesS \beta') \toS \boolS\) relational, with \(\castpath{\alpha}{\alpha'}\) and \(\castpath{\beta}{\beta'}\).
The encoder loosens \(S\) along \(\mathsf{graph}\) to a fresh \(\lvar{S} : (\alpha' \timesS \beta') \toS \boolS\), asserts \(\spec{\lvar{S}}\), and returns
\begin{equation}
  \bigl(\lamS \var{p} : \alpha' \timesS \beta' \cdot\; \appS{\lvar{S}}{\var{p}} \orS \appS{T}{\var{p}},\;\; (\alpha' \timesS \beta') \toS \boolS\bigr).
  \label{eq:beer-union-graph}
\end{equation}
On the running example, \(\var{f} : \intS \toS \optS\intS\) (flagged) and \(\var{g} : (\intS \timesS \intS) \toS \boolS\), so \(\var{f} \cupB \var{g}\) is encoded by \cref{eq:beer-union-graph}: \(\var{f}\) is loosened to the graph representation \(\lvar{f}\), and the union is the pointwise disjunction \(\lamS \var{p} \cdot \appS{\lvar{f}}{\var{p}} \orS \appS{\var{g}}{\var{p}}\).
Two further paths (\(\mathsf{fun}\), when both operands are functional, and \(\mathsf{chpred}\)) follow the same pattern at their respective common types; paths that cannot arise from B set operations (\(\mathsf{refl}\) on base types, \(\mathsf{opt}\), \(\mathsf{pair}\)) are rejected.

\paragraph{Equality.}
For \(A : \sigma_A\), \(B : \sigma_B\) ({\small\leancodeptr{castEq}{Encoder/Loosening/Loosening.html\#castEq}}):
if \(\sigma_A = \sigma_B\), emit \(A \eqS B\); if \(\castable{\sigma_A}{\sigma_B}\), loosen \(A\) and emit \((\lvar{A} \eqS B) \andS \spec{\lvar{A}}\); symmetrically if \(\castable{\sigma_B}{\sigma_A}\).
Here the specification travels \emph{inside} the returned formula rather than as a global assertion---a scoping necessity, since equalities may occur under binders.

\paragraph{Membership.}
For \(x : \sigma\) and a set operand \(S : \sigma' \toS \boolS\) ({\small\leancodeptr{castMembership}{Encoder/Loosening/Loosening.html\#castMembership}}), membership is predicate application, modulo loosening:
\[
  x \inB S \encto
  \begin{cases}
    (\appS{S}{x},\, \boolS) & \text{if } \sigma = \sigma', \\[2pt]
    (\spec{\lvar{x}} \andS \appS{S}{\lvar{x}},\, \boolS) & \text{if } \castable{\sigma}{\sigma'},\; \sigma \neq \sigma', \\[2pt]
    (\spec{\lvar{S}} \andS \appS{\lvar{S}}{x},\, \boolS) & \text{if } \castable{\sigma'}{\sigma},
  \end{cases}
\]
loosening the element up to the predicate's domain, or the predicate (along \(\mathsf{chpred}\)) up to the element's type.
When \(S\) is instead a flagged, functional operand \(S : \sigma'_1 \toS \optS\sigma'_2\) and \(x : \sigma_1 \timesS \sigma_2\) a pair, membership becomes a defined-application equation; in the representative subcase \(\castable{\sigma_1}{\sigma'_1}\), \(\castable{\sigma_2}{\sigma'_2}\):
\[
  x \inB S \encto \bigl(\spec{\lvar{x}} \andS (\appS{S}{(\fstS \lvar{x})} \eqS \someS (\sndS \lvar{x})),\;\boolS\bigr),
\]
and the three remaining subcases dispatch symmetrically, loosening \(\fstS x\), \(\sndS x\), or \(S\) itself componentwise.

\paragraph{Application.}
B application \(\appB f\ x\) must produce a \emph{value}, not a predicate, so a graph representation must be cast in the opposite, tightening direction---the one place where this is sound, because B's well-definedness (\Cref{sec:b-wd}) guarantees \(x \in \domB(f)\) and functionality of \(f\).
For \(f : (\tau \timesS \sigma) \toS \boolS\) and \(x : \xi\) with \(\castable{\tau}{\xi}\) ({\small\leancodeptr{castApp}{Encoder/Loosening/Loosening.html\#castApp}}), the encoder loosens \(f\) to \(\lvar{f} : (\xi \timesS \sigma) \toS \boolS\), introduces a fresh \emph{function witness} \(\var{f}^{!!} : \xi \toS \optS\sigma\) pinned to the graph by the assertion
\begin{equation}
  \forallS \var{u} : \xi,\, \var{v} : \sigma \cdot\;
  \appS{\lvar{f}}{\pairS{\var{u}}{\var{v}}} \eqS \bigl(\appS{\var{f}^{!!}}{\var{u}} \eqS \someS \var{v}\bigr),
  \label{eq:beer-castapp-spec}
\end{equation}
and returns
\(
  \bigl(\theS (\appS{\var{f}^{!!}}{x'}),\; \sigma\bigr)
\), where \(\theS\) extracts the payload of a defined option (\Cref{ch:smt}) and is unconstrained on \(\noneS\).
When \(f\) is already functional, \(f : \tau \toS \optS\sigma\), the witness is unnecessary and the encoder emits \(\theS(\appS{f}{x'})\) directly, loosening \(f\) or \(x\) as the types require.

\subsection{Binders and the universal quantifier}
\label{sec:beer-rules-quantifier}

The last rule is the bounded universal quantifier, whose translation must solve a scoping problem that no other rule has.
Encoding the body \(P\) and the domain guard may itself invoke the cast layer, producing fresh helper variables (\(\lvar{x}\), \(\var{f}^{!!}\), \dots) together with their specifications.
Helpers generated for the \emph{bound} variables \(\vec{\var{v}}\) cannot be declared globally: their values depend on the quantified variables and would otherwise escape their scope.
The encoder therefore re-scopes them \emph{inside} the quantifier, universally bound and guarded by their specifications.
With \(D \encto (D', \sigma \toS \boolS)\), fresh binders \(\vec{\var{z}}\) of types \(\vec{\tau}\), helpers \(\var{h}_1, \dots, \var{h}_k\) collected while encoding the guard and the body, and \(P'\) the encoded body with \(\vec{\var{v}}\) replaced by \(\vec{\var{z}}\):
\begin{equation}
  \forallB \vec{\var{v}} \inB D \cdot P
  \;\encto\;
  \Bigl(\forallS \vec{\var{z}} : \vec{\tau} \cdot\;
    \forallS \var{h}_1 \cdots \forallS \var{h}_k \cdot\;
    \spec{{\var{h}_1}} \impS \cdots \impS \spec{{\var{h}_k}} \impS
    \bigl( (\vec{\var{z}} \inB D)' \impS P' \bigr),\;\;
  \boolS\Bigr)
  \label{eq:beer-forall-shape}
\end{equation}
where \((\vec{\var{z}} \inB D)'\) abbreviates the membership encoding of \Cref{sec:beer-rules-cast}.
The binder types \(\vec{\tau}\) are the component types of \(\sigma\), except that a bound variable flagged as functional is given its tightened functional type, mirroring the treatment of declared variables.
The guarded-universal shape \(\forallS \var{h} \cdot \spec{\var{h}} \impS \cdots\)---rather than the seemingly natural \(\existsS \var{h} \cdot \spec{\var{h}} \andS \cdots\)---is a soundness matter, not a stylistic one: the argument, which hinges on the fact that specifications determine their helpers uniquely, is given with the soundness proof of this case (\Cref{sec:beer-case-forall}).

On the running example, the goal \(\var{f} \cupB \var{g} \subseteqB X \prodB Y\) is elaborated to \(\forallB \var{x} \inB \var{f} \cupB \var{g} \cdot \var{x} \inB X \prodB Y\), whose encoding instantiates \cref{eq:beer-forall-shape}: the domain encoding triggers the graph-path union \cref{eq:beer-union-graph}, whose loosened \(\lvar{f}\) becomes the helper \(\var{h}_1\) with \(\spec{{\var{h}_1}}\) the graph specification of \Cref{ch:loosening}, and the membership guard and body are encoded by the rules above.
The complete emitted script is shown in \Cref{sec:beer-proof-obligations}.
